\begin{table}[H]
    \centering
    \caption{Schema of the \texttt{header} table.}
    \label{tab:header-schema}
    \begin{tabular}{lll}
    \toprule
    \textbf{Column}                      & \textbf{Type} & \textbf{Description}                         \\ \midrule
    \texttt{id}             & INTEGER       & Identifier of the header record.             \\
    \texttt{created}        & INTEGER       & Creation time as a Unix timestamp.           \\
    \texttt{device\_id}     & TEXT          & Identifier of the device producing the data. \\
    \texttt{channels}       & INTEGER       & Number of spectrum channels.                 \\
    \texttt{calibration}    & TEXT          & Calibration coefficients as a json list.     \\
    \texttt{concat} &
      INTEGER &
      \begin{tabular}[c]{@{}l@{}}Indicates whether the channels have been concatenated \\ by a factor of 1 (not concatinated), 2, 4 or 8.\end{tabular} \\
    \texttt{save\_interval} & REAL          & Interval between data saves.                 \\
    \texttt{meta}           & TEXT          & Additional metadata as json.                 \\ \bottomrule
    \end{tabular}
\end{table}


\begin{table}[H]
    \centering
    \caption{Schema of the \texttt{summary} table.}
    \label{tab:summary-schema}
    \begin{tabular}{lll}
        \toprule
        \textbf{Column} & \textbf{Type} & \textbf{Description} \\
        \midrule
        \texttt{id}
            & INTEGER
            & Identifier of the summary record. \\

        \texttt{accumulated\_spectrum}
            & BLOB
            & \begin{tabular}[c]{@{}l@{}}Accumulated spectrum data, \\ zlib compressed array. \end{tabular} \\

        \texttt{total\_duration}
            & REAL
            & Total duration of the recording in [s]. \\

        \texttt{total\_dose}
            & REAL
            & Total accumulated dose in [\unit{\micro\sievert}]. \\

        \texttt{last\_update}
            & INTEGER
            & \begin{tabular}[c]{@{}l@{}}Timestamp of the last update as \\ a Unix timestamp. \end{tabular} \\

        \bottomrule
    \end{tabular}
\end{table}


\begin{table}[H]
    \centering
    \caption{Schema of the \texttt{spectrogram} table.}
    \label{tab:spectrogram-schema}
    \begin{tabular}{lll}
        \toprule
        \textbf{Column} & \textbf{Type} & \textbf{Description} \\
        \midrule
        \texttt{id}
            & INTEGER
            & Identifier of the spectrogram record. \\

        \texttt{timestamp}
            & REAL
            & Timestamp of the measurement as a Unix timestamp. \\

        \texttt{avg\_cps}
            & INTEGER
            & Average counts per second multiplied by 1000. \\

        \texttt{avg\_dr}
            & INTEGER
            & Average dose rate as [\unit{\micro\sievert\per\second}*1000]. \\

        \texttt{temperature}
            & INTEGER
            & Temperature during the measurement in [\unit{\celsius}*1000]. \\

        \texttt{latitude}
            & REAL
            & Latitude of the measurement location. \\

        \texttt{longitude}
            & REAL
            & Longitude of the measurement location. \\

        \texttt{spectrum}
            & BLOB
            & Recorded spectrum during interval, zlib compressed array. \\

        \texttt{meta}
            & BLOB
            & \begin{tabular}[c]{@{}l@{}}Additional metadata as a json string \\ encoded with utf-8 and zlib compressed.\end{tabular} \\

        \bottomrule
    \end{tabular}
\end{table}
